% unit U4
\section{The Full \treeK Target}\label{sec:target}
% unit macros
\providecommand{\yields}{\mathrel{\Longrightarrow}}
\mathtgt

This section gives the whole target of \treeK. It adds to the fragment of Section~\ref{sec:mini} what \Rone adds to its own, now with capture sets, and three things more: capture names with capture members, a store in the theorem about capture evidence, and boxes. Figures~\ref{fig:target-syntax} to~\ref{fig:target-capev} give the calculus where it differs from report~1's Figures~\ref{R1-fig:target-syntax} to~\ref{R1-fig:machine}. Its theorems are proven in \treeK, and Table~\ref{tab:matrix} lists those that another tree lacks or states differently.

\paragraph{A program with a capability member.}
The reactor of Section~\ref{sec:cc-later} lets a handler of a \lstinline{GUIReactor} use \lstinline{ui} and \lstinline{log} and nothing else. The Lean's example C2 has this shape over two platform capabilities $\kappa_1$ and $\kappa_2$. Its object type, the translation of a DOT type that Section~\ref{sec:source} writes out, is
\[
\begin{array}{@{}l@{}}
\mathit{Obj} = \mu(\self.\,[\,\plec{\cempty}{\cset{\sel{\self}{C}}},\ \plec{\cset{\sel{\self}{C}}}{\cset{\kappa_1,\kappa_2}},\ \phas{\mathit{run}}, \\
\qquad\qquad \ple{\sel{\self}{\mathit{run}}}{\pit[w]{\Unit}{\Unit}},\ \plec{\cset{\sel{\self}{\mathit{run}}}}{\cset{\sel{\self}{C}}}\,]) .
\end{array}
\]
The promise is that a call of $\mathit{run}$ is charged to the member's bound $\cset{\kappa_1,\kappa_2}$. Section~\ref{sec:target-members} gives the evidence for that charge, and Section~\ref{sec:target-canon} the theorem that makes it mean something.

\begin{figure}[t]
\small
\color{tgtcol}
\[
\begin{array}{@{}lrcl@{}}
\text{Capture atoms} & c & ::= & x \mid \kappa \mid \sel{x}{\ell} \\
\text{Capture sets} & A, C, D, U & ::= & \cset{c_1, \ldots, c_n} \\
\text{Shapes} & S & ::= & \bot \mid \sel{x}{\ell} \mid \pit[x]{T}{T'} \mid \obj[x]{\tel} \mid \boxt{T} \\
\text{Types} & T & ::= & \capt{S}{C} \\
\text{Propositions} & P & ::= & \ple{S}{S'} \mid \peq{S}{S'} \mid \phas{\ell} \mid \pbnd{S} \mid \plec{C}{D} \mid \peqc{C}{D} \\[2pt]
\text{Shape inclusion} & e & ::= & \text{report~1's inclusion evidence, at shapes} \mid \ev{boxed}\,d \\
\text{Type inclusion} & d & ::= & \ev{capt}\,e\,f \\
\text{Capture inclusion} & f, g & ::= & \ev{refl}\,C \mid \ev{trans}\,f\,g \mid \ev{elem}\,C\,D \mid \ev{union}\,f\,g \\
& & \mid & \ev{capvar}\,a \mid \ev{member}\,a\,e\,i \mid \ev{eqToLe}\,\psi \\
\text{Capture equality} & \psi & ::= & \ev{refl}\,C \mid \ev{symm}\,\psi \mid \ev{trans}\,\psi\,\psi \mid \ev{defC}\,x\,\ell \mid \ev{member}\,a\,e\,i \\
\text{Capture holes} & H & ::= & \ev{leC}\,j \mid \ev{eqC}\,j \mid \ev{eqSymC}\,j \\
\text{Capture steps} & \theta & ::= & \ev{closed}\,f \mid \ev{incl}\,C\,D \\
\text{Capture sides} & p & ::= & [\,] \mid \theta :: p \\
\text{Templates} & \tau & ::= & \ev{le}\,q\,H\,q' \mid \ev{eq}\,j\,\beta \mid \ev{has}\,j \mid \ev{bnd}\,e \mid \ev{leC}\,p\,H\,p' \mid \ev{eqC}\,j\,\beta \\[2pt]
\text{Atoms} & a, b & ::= & x \mid \ev{cast}\,a\,d \mid \ev{fold}\,\tel\,a \mid \ev{unfold}\,a \mid \ev{both}\,\tel\,\tel\,a\,a \mid \ev{recap}\,a\,f \\
\text{Terms} & t, u & ::= & a \mid v \mid a\,b \mid \ev{proj}\,a\,\ell\,h \mid \letc{x}{t}{u}{U}{f} \mid \ev{cast}\,t\,d \mid \unboxc{a}{U}{f} \\
\text{Values} & v & ::= & \lamc{A}{x}{T}{t}{g} \mid \litc{A}{W}{W^{c}}{F} \mid \ev{box}\,a \mid \ev{cast}\,v\,d \\
\text{Witnesses} & W & ::= & \tnil \mid W, \ell = S \\
\text{Capture witnesses} & W^{c} & ::= & \tnil \mid W^{c}, \ell = C \\
\text{Fields} & F & ::= & \tnil \mid F, \ell = t\,\ann{g} \\[2pt]
\text{Capture bounds} & B & ::= & \cstar \mid \scopeb \mid \mathsf{upper}\,C \mid \instb{C} \\
\text{Contexts} & \Gamma & ::= & \cdot \mid \Gamma, x : T \mid \Gamma, x : T \mathrel{:=} \langle W, W^{c}, \vec{\ell}\,\rangle \mid \Gamma, \kappa : B \\
\text{Stores} & \sigma & ::= & \cdot \mid \sigma, x = v \mid \sigma, \kappa : B \\
\text{Continuations} & K & ::= & \cdot \mid K \triangleright \frlet{x}{u}{U}{f} \mid K \triangleright \ev{cast}\,d
\end{array}
\]
\[
\begin{array}[t]{@{}lcl@{}}
\rt{\ev{recap}\,a\,f} & = & \rt{a} \\
\uses{a} & = & \cset{\rt{a}} \\
\uses{a\,b} & = & \cset{\rt{a}, \rt{b}} \\
\uses{\ev{proj}\,a\,\ell\,h} & = & \cset{\rt{a}} \\
\uses{\ev{cast}\,t\,d} & = & \uses{t} \\
\uses{\unboxc{a}{U}{f}} & = & \cset{\rt{a}} \cup U \\
\uses{K \triangleright \ev{cast}\,d} & = & \uses{K}
\end{array}
\quad
\begin{array}[t]{@{}lcl@{}}
\insp{a\,b} & = & \rt{a} \\
\insp{\ev{proj}\,a\,\ell\,h} & = & \rt{a} \\
\insp{\unboxc{a}{U}{f}} & = & \rt{a} \\
\insp{t} & & \text{undefined otherwise} \\
\annot{\litc{A}{W}{W^{c}}{F}} & = & A \\
\annot{\ev{box}\,a} & = & \cempty \\
\annot{\ev{cast}\,v\,d} & = & \annot{v}
\end{array}
\]
\color{black}
\caption{The \treeK target: what it adds to the fragment of Figure~\ref{fig:mini-syntax}. $\top$ abbreviates $\obj{\tnil}$. A capture name $\sel{x}{\ell}$ stands in a capture set, a block name $\sel{x}{\ell}$ where a shape does. In $\ev{le}\,q\,H\,q'$ the sides and the hole are report~1's, in $\ev{leC}\,p\,H\,p'$ capture sides and a capture hole. The flag $\beta$ of $\ev{eq}\,j\,\beta$ and $\ev{eqC}\,j\,\beta$ is $\mathit{true}$ or $\mathit{false}$, and $\mathit{true}$ flips the equality. A binder $\kappa : (\instb{C})$ is written $\kappa \instb{C}$. Binders are as in report~1's Figure~\ref{R1-fig:target-syntax}, and a literal binds $\self$ in $W$, $W^{c}$ and $F$. $\mathit{root}$ is report~1's with the clause for $\ev{recap}$ added. The other equations shown are those of Figure~\ref{fig:mini-syntax} that change and the new ones. The rest carry over.}
\label{fig:target-syntax}
\end{figure}

\subsection{What the full target adds}\label{sec:target-delta}

A type is a shape with a capture set, $\capt{S}{C}$, and report~1's types are the shapes, with the box $\boxt{T}$ of Section~\ref{sec:target-boxes} added. Report~1's inclusion evidence becomes shape inclusion $e$, and a type inclusion $\ev{capt}\,e\,f$ pairs it with capture inclusion $f$. The function shape takes type inclusions for its domain and codomain (\rn{Sh-Pi}), so a cast checks the capture sets on both sides of an arrow. Casts supply the subsumption that the fragment did without. A capture set may contain a capture name $\sel{x}{\ell}$, the capture set named $\ell$ in the block of $x$, and telescopes gain capture inclusions and equalities.

Atoms gain $\ev{recap}\,a\,f$, which keeps the shape of $a$ and takes any capture set that $f$ puts above $\cset{\rt{a}}$ (\rn{A-Recap}). With $f = \ev{refl}\,\cset{x}$ it types $x$ at $\capt{S}{\cset{x}}$, the variable rule of capture calculi (Section~\ref{sec:source}). Every wrapper keeps the root, so \rn{C-Var} and \rn{C-Member} read the same variable however an atom is wrapped (\lean{C4\_capvar}, \lean{C4\_member}, decided). A projection has the block name of the field at its capture name (\rn{T-Proj}). A literal carries an annotation $A$ and capture witnesses $W^{c}$, and each of its fields carries closing evidence against $A$ and the self (\rn{F-Cons}). Figure~\ref{fig:target-rules} gives the rules that are new or changed. Its steps read two head forms that report~1's Figure~\ref{R1-fig:forms} lacks: $\form{eqv}\,\varphi$ for closed equality evidence, which report~1 folds into $\form{id}$, and $\form{boxed}\,d$ for an inclusion of boxes (Section~\ref{sec:mech-trees}).

No rule for terms adds a capture binder to a context, and no step allocates one. So no term abstracts over a capture set, and DOT quantifies over capture sets through members.

\begin{figure}[t]
\small
\color{tgtcol}
\[
\begin{array}{@{}l@{\quad}l@{}}
\rn{L-Capt} & \shpj{\Gamma}{e}{S}{S'} \quad \clej{\Gamma}{f}{C}{C'} \yields \lej{\Gamma}{\ev{capt}\,e\,f}{\capt{S}{C}}{\capt{S'}{C'}} \\
\rn{Sh-Pi} & \lej{\Gamma}{d}{T_2}{T_1} \quad \lej{\Gamma, x : T_2}{d'}{T'_1}{T'_2} \yields \shpj{\Gamma}{\ev{pi}\,d\,d'}{\pit{T_1}{T'_1}}{\pit{T_2}{T'_2}} \\
\rn{Sh-Boxed} & \lej{\Gamma}{d}{T}{T'} \yields \shpj{\Gamma}{\ev{boxed}\,d}{\boxt{T}}{\boxt{T'}} \\[3pt]
\rn{A-Unfold} & \atomj{\Gamma}{a}{\capt{(\obj{\tel})}{C}} \yields \atomj{\Gamma}{\ev{unfold}\,a}{\capt{(\obj{\weak{(\tel[\self \mapsto \rt{a}])}})}{C}} \\
\rn{A-Fold} & \atomj{\Gamma}{a}{\capt{(\obj{\weak{(\tel[\self \mapsto \rt{a}])}})}{C}} \yields \atomj{\Gamma}{\ev{fold}\,\tel\,a}{\capt{(\obj{\tel})}{C}} \\
\rn{A-Both} & \atomj{\Gamma}{a}{\capt{(\obj{\tel_1})}{C}} \quad \atomj{\Gamma}{b}{\capt{(\obj{\tel_2})}{C}} \quad \rt{b} = \rt{a} \\
& \qquad \yields \atomj{\Gamma}{\ev{both}\,\tel_1\,\tel_2\,a\,b}{\capt{(\obj{\tel_1 \mathbin{+\!\!+} \tel_2})}{C}} \\
\rn{A-Recap} & \atomj{\Gamma}{a}{\capt{S}{C}} \quad \clej{\Gamma}{f}{\cset{\rt{a}}}{C'} \yields \atomj{\Gamma}{\ev{recap}\,a\,f}{\capt{S}{C'}} \\[3pt]
\rn{T-App} & \atomj{\Gamma}{a}{\capt{(\pit[z]{T}{T'})}{C}} \quad \atomj{\Gamma}{b}{T} \yields \tmj{\Gamma}{a\,b}{T'[z \mapsto \rt{b}]} \\
\rn{T-Proj} & \atomj{\Gamma}{a}{T} \quad \hasj{\Gamma}{h}{\rt{a}}{\ell} \yields \tmj{\Gamma}{\ev{proj}\,a\,\ell\,h}{\capt{(\sel{\rt{a}}{\ell})}{\cset{\sel{\rt{a}}{\ell}}}} \\
\rn{T-Unbox} & \atomj{\Gamma}{a}{\capt{(\boxt{\capt{S}{C}})}{D}} \quad \clej{\Gamma}{f}{C}{U} \yields \tmj{\Gamma}{\unboxc{a}{U}{f}}{\capt{S}{C}} \\[3pt]
\rn{V-Obj} & \Gamma, \self : \capt{(\obj{\ofLit{W, W^{c}}{\dom F}})}{A} \mathrel{:=} \langle W, W^{c}, \dom F\rangle \vdash^{A}_{\!f} F \\
& \qquad \yields \valj{\Gamma}{\litc{A}{W}{W^{c}}{F}}{\capt{(\obj{\ofLit{W, W^{c}}{\dom F}})}{A}} \\
\rn{V-Box} & \atomj{\Gamma}{a}{T} \yields \valj{\Gamma}{\ev{box}\,a}{\capt{(\boxt{T})}{\cempty}} \\
\rn{F-Cons} & \Gamma \vdash^{A}_{\!f} F \quad \tmj{\Gamma}{t}{\capt{(\sel{\self}{\ell})}{\cset{\sel{\self}{\ell}}}} \quad \clej{\Gamma}{g}{\uses{t}}{\weak{A} \cup \cset{\self}} \yields \Gamma \vdash^{A}_{\!f} F, \ell = t\,\ann{g} \\[3pt]
\rn{S-Val} & \storej{\sigma}{\Gamma} \quad v \text{ is not a cast} \quad \valj{\Gamma}{v}{T} \yields \storej{\sigma, x = v}{\Gamma, x : T \mathrel{:=} \langle \mathit{wit}(v), \mathit{wit}^{c}(v), \mathit{fields}(v)\rangle} \\
\rn{S-Cap} & \storej{\sigma}{\Gamma} \yields \storej{\sigma, \kappa : B}{\Gamma, \kappa : B} \\
\rn{K-Cast} & \lej{\Gamma}{d}{T}{T'} \quad \contj{\Gamma}{K}{T'}{T''} \yields \contj{\Gamma}{K \triangleright \ev{cast}\,d}{T}{T''}
\end{array}
\]
\[
\begin{array}{@{}l@{\quad}l@{\ }c@{\ }l@{}}
\rn{E-Alloc} & \mst{\sigma}{K \triangleright \frlet{x}{u}{U}{f}}{v} &\step& \mst{(\sigma, x = \mathit{core}(v))}{\weak{K}}{\mathit{adjust}(u, v)} \\
\rn{E-Rename} & \mst{\sigma}{K \triangleright \frlet{x}{u}{U}{f}}{a} &\step& \mst{\sigma}{K}{u[x \mapsto a]} \\
\rn{E-AppCastRefl} & \mst{\sigma}{K}{a\,b} &\step& \mst{\sigma}{K}{t_0[z \mapsto b]} \\
& \multicolumn{3}{l}{\quad \text{if } \sigma(\rt{a}) = \lamc{A}{z}{T_0}{t_0}{g},\ a \text{ is not a variable},\ \chn{\sigma}{a}{n}{\form{id}} \text{ or } \chn{\sigma}{a}{n}{\form{eqv}\,\varphi}} \\
\rn{E-UnboxRefl} & \mst{\sigma}{K}{\unboxc{a}{U}{f}} &\step& \mst{\sigma}{K}{b} \\
& \multicolumn{3}{l}{\quad \text{if } \sigma(\rt{a}) = \ev{box}\,b,\ \chn{\sigma}{a}{n}{\form{id}} \text{ or } \chn{\sigma}{a}{n}{\form{eqv}\,\varphi}} \\
\rn{E-UnboxCast} & \mst{\sigma}{K}{\unboxc{a}{U}{f}} &\step& \mst{\sigma}{K}{\ev{cast}\,b\,d} \quad \text{if } \sigma(\rt{a}) = \ev{box}\,b,\ \chn{\sigma}{a}{n}{\form{boxed}\,d}
\end{array}
\]
\color{black}
\caption{The rules and steps of \treeK that are new or changed. \rn{T-Let}, \rn{V-Lam}, \rn{S-Nil}, \rn{K-Nil}, \rn{K-Let}, \rn{St}, \rn{E-Let}, \rn{Run-Refl} and \rn{Run-Step} are those of Figures~\ref{fig:mini-typing} and~\ref{fig:mini-machine}. The other rules and steps are those of report~1's Figures~\ref{R1-fig:evidence} to~\ref{R1-fig:machine} at shapes, without the rules for $\ev{pi}$ and $\ev{leTrans}$, and an atom premise of $\ev{member}$ takes any capture set. Report~1's rules for variable and cast atoms are \rn{A-Var} and \rn{A-Cast} here. Its other steps are \rn{E-CastPush}, \rn{E-CastVal}, \rn{E-CastAtom}, \rn{E-App} through a variable, \rn{E-AppCast} through a chain $\form{pi}\,d\,c$, and \rn{E-Proj}. $\mathit{core}$ and $\mathit{adjust}$ are those of its Figure~\ref{R1-fig:mini-machine}. Substituting an atom, as in $u[x \mapsto a]$, puts $a$ where an atom stands and $\rt{a}$ in types and capture sets, as in report~1. $\vdash^{A}_{\!f}$ types the fields of a literal annotated $A$. $\mathit{wit}(v)$, $\mathit{wit}^{c}(v)$ and $\mathit{fields}(v)$ are the witnesses, capture witnesses and field labels of $v$, empty for a lambda and a box.}
\label{fig:target-rules}
\end{figure}

\subsection{Capture names and capture members}\label{sec:target-members}

A source capture member becomes two capture inclusions on the capture name of the self, as in $\mathit{Obj}$. A field gets a third proposition after report~1's two, its declared capture set. Figure~\ref{fig:source-trans} gives both clauses. A literal defines its capture names by its capture witnesses. Its precise telescope $\ofLit{W, W^{c}}{\vec{\ell}}$ lists one equality $\peq{\sel{\self}{\ell}}{W(\ell)}$ per witness, then one capture equality $\peqc{\cset{\sel{\self}{\ell}}}{W^{c}(\ell)}$ per capture witness, then one presence per field. A store binder records the capture witnesses, and $\Gamma(x) \ni \ell \mathrel{:=_{c}} C$ says that the witness of $x$ at $\ell$ is $C$, with $C = \cempty$ for an unlisted $\ell$.

Figure~\ref{fig:target-capev} gives capture evidence in full. \rn{C-Member} eliminates a capture inclusion at an atom, as report~1's $\ev{member}$ eliminates an inclusion, and \rn{CE-Def} reads a capture name of a transparent binder. The capture template $\ev{leC}\,p\,H\,p'$ proves a target capture inclusion from the source capture proposition that its hole names. Its sides chain closed capture evidence, weakened under the self, and syntactic inclusions $\ev{incl}\,X\,Y$, which may mention the self.

C2's client shows the charge (\lean{C2\_client}, decided). In the context $\kappa_1 : \cstar, \kappa_2 : \cstar, x : \capt{\mathit{Obj}}{\cset{\kappa_1,\kappa_2}}, u : \Unit$ it reads $\mathit{run}$ off $x$ and calls it:
\[
\begin{array}{@{}l@{}}
\klet r = \ev{cast}\,(\ev{proj}\,x\,\mathit{run}\,h_2)\,(\ev{capt}\,e_3\,f_4) \kin r\,u \\
\qquad \ann{\cset{\kappa_1,\kappa_2};\ \ev{union}\,(\ev{trans}\,(\ev{capvar}\,r)\,f_1)\,(\ev{trans}\,(\ev{capvar}\,u)\,(\ev{elem}\,\cempty\,\cset{\kappa_1,\kappa_2}))} .
\end{array}
\]
Here $r$ has the type $\capt{(\pit[w]{\Unit}{\Unit})}{\cset{\sel{x}{C}}}$, since $h_2$, $e_3$, $f_4$ and $f_1$ are $\ev{member}\,x\,(\ev{refl}\,\mathit{Obj})\,i$ at $i = 2, 3, 4, 1$. The avoidance evidence charges the call. \rn{C-Var} puts $\cset{r}$ below $\cset{\sel{x}{C}}$, and $f_1$ reads the member's upper bound at $x$, $\cset{\sel{x}{C}} \le \cset{\kappa_1,\kappa_2}$. The client is checked once, for every object of type $\mathit{Obj}$.

\begin{figure}[t]
\small
\color{tgtcol}
\[
\begin{array}{@{}l@{\quad}l@{}}
\rn{C-Refl} & \yields \clej{\Gamma}{\ev{refl}\,C}{C}{C} \\
\rn{C-Trans} & \clej{\Gamma}{f}{C_1}{C_2} \quad \clej{\Gamma}{g}{C_2}{C_3} \yields \clej{\Gamma}{\ev{trans}\,f\,g}{C_1}{C_3} \\
\rn{C-Elem} & C \subseteq D \yields \clej{\Gamma}{\ev{elem}\,C\,D}{C}{D} \\
\rn{C-Union} & \clej{\Gamma}{f}{C_1}{D} \quad \clej{\Gamma}{g}{C_2}{D} \yields \clej{\Gamma}{\ev{union}\,f\,g}{C_1 \cup C_2}{D} \\
\rn{C-Var} & \atomj{\Gamma}{a}{\capt{S}{C}} \yields \clej{\Gamma}{\ev{capvar}\,a}{\cset{\rt{a}}}{C} \\
\rn{C-Member} & \atomj{\Gamma}{a}{\capt{S}{D}} \quad \shpj{\Gamma}{e}{S}{\obj{\tel}} \quad \tat{\tel}{i}{\plec{C_1}{C_2}} \\
& \qquad \yields \clej{\Gamma}{\ev{member}\,a\,e\,i}{C_1[\self \mapsto \rt{a}]}{C_2[\self \mapsto \rt{a}]} \\
\rn{C-EqToLe} & \ceqj{\Gamma}{\psi}{C}{D} \yields \clej{\Gamma}{\ev{eqToLe}\,\psi}{C}{D} \\[3pt]
\rn{CE-Refl} & \yields \ceqj{\Gamma}{\ev{refl}\,C}{C}{C} \\
\rn{CE-Symm} & \ceqj{\Gamma}{\psi}{C}{D} \yields \ceqj{\Gamma}{\ev{symm}\,\psi}{D}{C} \\
\rn{CE-Trans} & \ceqj{\Gamma}{\psi}{C_1}{C_2} \quad \ceqj{\Gamma}{\psi'}{C_2}{C_3} \yields \ceqj{\Gamma}{\ev{trans}\,\psi\,\psi'}{C_1}{C_3} \\
\rn{CE-Def} & \Gamma(x) \ni \ell \mathrel{:=_{c}} C \yields \ceqj{\Gamma}{\ev{defC}\,x\,\ell}{\cset{\sel{x}{\ell}}}{C} \\
\rn{CE-Member} & \atomj{\Gamma}{a}{\capt{S}{D}} \quad \shpj{\Gamma}{e}{S}{\obj{\tel}} \quad \tat{\tel}{i}{\peqc{C_1}{C_2}} \\
& \qquad \yields \ceqj{\Gamma}{\ev{member}\,a\,e\,i}{C_1[\self \mapsto \rt{a}]}{C_2[\self \mapsto \rt{a}]} \\[3pt]
\rn{H-LeC} & \tat{\tel}{j}{\plec{X}{Y}} \yields \tel \vdash \ev{leC}\,j : \plec{X}{Y} \\
\rn{H-EqC} & \tat{\tel}{j}{\peqc{X}{Y}} \yields \tel \vdash \ev{eqC}\,j : \plec{X}{Y} \\
\rn{H-EqSymC} & \tat{\tel}{j}{\peqc{Y}{X}} \yields \tel \vdash \ev{eqSymC}\,j : \plec{X}{Y} \\[3pt]
\rn{Sd-Closed} & \clej{\Gamma}{f}{C}{D} \yields \Gamma \vdash \ev{closed}\,f : \weak{C} \le \weak{D} \\
\rn{Sd-Incl} & X \subseteq Y \yields \Gamma \vdash \ev{incl}\,X\,Y : X \le Y \\
\rn{Sd-Nil} & \yields \Gamma \vdash [\,] : X \le X \\
\rn{Sd-Cons} & \Gamma \vdash \theta : X \le Y \quad \Gamma \vdash p : Y \le Z \yields \Gamma \vdash \theta :: p : X \le Z \\[3pt]
\rn{M-LeC} & \tel \vdash H : \plec{X}{Y} \quad \Gamma \vdash p : D_1 \le X \quad \Gamma \vdash p' : Y \le D_2 \yields \morj{\Gamma}{\ev{leC}\,p\,H\,p'}{\tel}{\plec{D_1}{D_2}} \\
\rn{M-EqC} & \tat{\tel}{j}{\peqc{C_1}{C_2}} \yields \morj{\Gamma}{\ev{eqC}\,j\,\mathit{false}}{\tel}{\peqc{C_1}{C_2}} \\
\rn{M-EqSymC} & \tat{\tel}{j}{\peqc{C_1}{C_2}} \yields \morj{\Gamma}{\ev{eqC}\,j\,\mathit{true}}{\tel}{\peqc{C_2}{C_1}}
\end{array}
\]
\[
\begin{array}{@{}lcl@{}}
\rootsn{\Gamma}{n}{\cempty} & = & \cempty \\
\rootsn{\Gamma}{n}{\cset{c} \cup C} & = & \rootsn{\Gamma}{n}{c} \cup \rootsn{\Gamma}{n}{C} \\
\rootsn{\Gamma}{n}{x} & = & \rootsn{\Gamma'}{n}{C} \quad \text{if } \Gamma = \Gamma', x : \capt{S}{C}, \Gamma'' \text{, opaque or transparent} \\
\rootsn{\Gamma}{n}{\kappa} & = & \cset{\kappa} \quad \text{if } \Gamma = \Gamma', \kappa : \cstar, \Gamma'' \\
\rootsn{\Gamma}{n}{\kappa} & = & \cset{\kappa} \quad \text{if } \Gamma = \Gamma', \kappa : \scopeb, \Gamma'' \\
\rootsn{\Gamma}{n}{\kappa} & = & \rootsn{\Gamma'}{n}{C} \quad \text{if } \Gamma = \Gamma', \kappa : \mathsf{upper}\,C, \Gamma'' \\
\rootsn{\Gamma}{n}{\kappa} & = & \rootsn{\Gamma'}{n}{C} \quad \text{if } \Gamma = \Gamma', \kappa \instb{C}, \Gamma'' \\
\rootsn{\Gamma}{0}{\sel{x}{\ell}} & = & \cempty \\
\rootsn{\Gamma}{n+1}{\sel{x}{\ell}} & = & \rootsn{\Gamma}{n}{C} \quad \text{if } \Gamma(x) \ni \ell \mathrel{:=_{c}} C \\
\rootsn{\Gamma}{n+1}{\sel{x}{\ell}} & = & \cempty \quad \text{if } x \text{ is opaque in } \Gamma \\[3pt]
c \in \roots{\Gamma}{C} & :\Leftrightarrow & c \in \rootsn{\Gamma}{n}{C} \text{ for some } n
\end{array}
\]
\color{black}
\caption{Capture evidence in full, with the capture templates of object coercions and the roots of a capture set with names and fuel. $\preccurlyeq_\Gamma$ and $\approx_\Gamma$ are defined from roots as in Figure~\ref{fig:mini-syntax}. The first equation that applies is used.}
\label{fig:target-capev}
\end{figure}

\subsection{Roots with names}\label{sec:target-roots}

The roots at the end of Figure~\ref{fig:target-capev} extend those of Figure~\ref{fig:mini-syntax} in two ways. First, a capture name follows the capture witness of its binder and spends one unit of fuel. A witness may name a label of its own block, so names can form a cycle. A cycle contributes no root, as report~1's cyclic alias resolves to $\top$. A capture name of an opaque binder has no roots, but over a store every term binder is transparent.

Second, a capture binder bound $\cstar$ or $\scopeb$ is its own root, and one bound $\mathsf{upper}\,C$ or $\instb{C}$ has the roots of $C$. These bounds are the scope, bounded and instance binders of \treeCC (Section~\ref{sec:levels-scopes}), and all three trees have them. No clause returns a term variable or a capture name, so every root is a capture binder bound $\cstar$ or $\scopeb$. \treeK constructs only $\cstar$, for the platform, but \rn{S-Cap} accepts any bound, and the theorems below hold over stores with any bounds.

\subsection{Canonical forms of closed capture evidence}\label{sec:target-canon}

Report~1's Theorem~\ref{R1-thm:canon} assumes a typed store, $\storej{\sigma}{\Gamma}$, and normalizes closed evidence to a typed head form. The normalizer drops the capture half of $\ev{capt}\,e\,f$, and the canonical form of capture evidence is a fact about roots, as in Theorem~\ref{thm:mini-cap}. That theorem held in every context. With capture members it cannot, since under a binder of type $\obj{[\,\plec{\cset{\kappa}}{\cempty}\,]}$ one $\ev{member}$ proves $\cset{\kappa} \le \cempty$, and $\kappa$ is a root of $\cset{\kappa}$ only (Example~\ref{ex:target-bad}). Over a typed store it holds again.

The member cases read a view of an atom, which has one proposition form per proposition of its type (report~1's Figure~\ref{R1-fig:forms}). At a capture proposition the form carries no data, and its typedness is the proposition at the root of the atom, read as a fact about the roots of capture sets. We call such a form a capture entry. A capture entry gets its fact from the store. A stored literal has one capture entry per capture witness, with the fact of \rn{CE-Def}, and a capture template composes its sides with the source capture entry.

\begin{theorem}[Canonical forms of closed capture evidence, \leanin{cap\_canon}{\treeK}]\label{thm:target-canon}
Let $\storej{\sigma}{\Gamma}$.
\begin{enumerate}
\item If $\clej{\Gamma}{f}{C}{D}$ then $\caple{\Gamma}{C}{D}$.
\item If $\ceqj{\Gamma}{\psi}{C}{D}$ then $\capeqv{\Gamma}{C}{D}$ (\lean{capeq\_canon}).
\item If $\atomj{\Gamma}{a}{\capt{S}{C}}$ then $\caple{\Gamma}{\cset{\rt{a}}}{C}$ (the capture conjunct of \lean{atom\_canon}).
\item If $\lej{\Gamma}{d}{\capt{S}{C}}{\capt{S'}{C'}}$ then $\caple{\Gamma}{C}{C'}$ (\lean{le\_canon\_cap}).
\end{enumerate}
\end{theorem}

\begin{proof}
By one mutual induction on typing derivations, together with report~1's Theorem~\ref{R1-thm:canon}, which the same induction proves at shapes (\lean{shape\_canon} and the view conjunct of \lean{atom\_canon}).
\begin{itemize}
\item \rn{C-Refl}, \rn{C-Trans}, \rn{C-Elem} and \rn{C-Union} as in Theorem~\ref{thm:mini-cap}. \rn{C-Var} is item 3, and \rn{C-EqToLe} is item 2.
\item \rn{C-Member} and \rn{CE-Member}, the only cases that read the store. Report~1's theorem gives the view of $a$ and the form of $e$, and pushes the view through the form to a typed view of the target telescope. Its entry at $i$ is a capture entry, and its typedness is the conclusion, $\caple{\Gamma}{C_1[\self \mapsto \rt{a}]}{C_2[\self \mapsto \rt{a}]}$, or $\approx_\Gamma$ for an equality.
\item \rn{CE-Refl}, \rn{CE-Symm} and \rn{CE-Trans}: $\approx_\Gamma$ is an equivalence. \rn{CE-Def}: a defined capture name has the roots of its witness, by the name clause (\lean{Ctx.Root\_name}).
\item Item 3. \rn{A-Var} is Lemma~\ref{lem:mini-var}, now at every fuel (\lean{Ctx.Root\_var}). \rn{A-Cast} is items 3 and 4, \rn{A-Fold}, \rn{A-Unfold} and \rn{A-Both} keep the capture set, and \rn{A-Recap} is item 1. Item 4. \rn{L-Capt} is item 1 at its capture half. \qedhere
\end{itemize}
\end{proof}

\begin{example}[A capture member through a literal, on paper]\label{ex:target-member}
Let the store hold $\kappa_1 : \cstar$, $\kappa_2 : \cstar$ and, at $y$, the literal $\litc{\cempty}{\tnil}{C = \cset{\kappa_1}}{\tnil}$. Its context $\Gamma$ types $y$ at $\capt{(\obj{\tel_y})}{\cempty}$ with $\tel_y = [\,\peqc{\cset{\sel{\self}{C}}}{\cset{\kappa_1}}\,]$ and records $\Gamma(y) \ni C \mathrel{:=_{c}} \cset{\kappa_1}$. Let $\mathit{Abs}$, the translation of the capture member of C2's object type, and the templates $m$ be
\[
\begin{array}{@{}l@{}}
\mathit{Abs} = \obj{[\,\plec{\cempty}{\cset{\sel{\self}{C}}},\ \plec{\cset{\sel{\self}{C}}}{\cset{\kappa_1,\kappa_2}}\,]} , \\
m = [\,\ev{leC}\,(\ev{incl}\,\cempty\,\cset{\kappa_1} :: [\,])\,(\ev{eqSymC}\,0)\,[\,],\ \ev{leC}\,[\,]\,(\ev{eqC}\,0)\,(\ev{incl}\,\cset{\kappa_1}\,\cset{\kappa_1,\kappa_2} :: [\,])\,] .
\end{array}
\]
So $a = \ev{cast}\,y\,(\ev{capt}\,(\ev{obj}\,\tel_y\,m)\,(\ev{refl}\,\cempty))$ has type $\capt{\mathit{Abs}}{\cempty}$, and $\ev{member}\,a\,(\ev{refl}\,\mathit{Abs})\,1 : \cset{\sel{y}{C}} \le \cset{\kappa_1,\kappa_2}$ reads the upper bound at $y$.

In the member case of item 1, the fact $\capeqv{\Gamma}{\cset{\sel{y}{C}}}{\cset{\kappa_1}}$ of the capture witness (\lean{Ctx.Root\_name}) passes through the second template of $m$ and its inclusion $\cset{\kappa_1} \subseteq \cset{\kappa_1,\kappa_2}$, so $\caple{\Gamma}{\cset{\sel{y}{C}}}{\cset{\kappa_1,\kappa_2}}$. Report~1's Example~\ref{R1-ex:mini-sel} makes the same computation. \lean{C5\_packing} in \treeK, built by hand from two decided parts, types a cast of this kind for a literal with the capture witness $\cset{\kappa_1}$, and no Lean statement computes the view.
\end{example}

\begin{corollary}[No closed escape of a star binder, \leanin{Store.Typed.no\_cap\_escape}{\treeK}]\label{cor:target-star}
Let $\storej{\sigma}{\Gamma}$ and $\Gamma(\kappa) = \cstar$. If $\kappa \notin \roots{\Gamma}{D}$ then no $f$ has $\clej{\Gamma}{f}{\cset{\kappa}}{D}$. In particular no $f$ has $\clej{\Gamma}{f}{\cset{\kappa}}{\cempty}$ (\lean{Store.Typed.no\_cap\_star\_le\_nil}).
\end{corollary}

\begin{proof}
$\kappa$ is a root of $\cset{\kappa}$, so item 1 of Theorem~\ref{thm:target-canon} would make it a root of $D$.
\end{proof}

\begin{example}[Bad capture bounds, revisited]\label{ex:target-bad}
The target typing of C3, the closure of Section~\ref{sec:problem-bad}, is over one star binder $\kappa : \cstar$. Its outer binder has the object type $\mathit{Bad} = \obj{[\,\plec{\cset{\kappa}}{\cempty}\,]}$, whose one proposition states the bad inclusion directly:
\[
\begin{array}{@{}l@{}}
\lambda^{\cempty}(x : \capt{\mathit{Bad}}{\cempty}).\,\lambda^{\cempty}(z : \capt{(\boxt{\capt{\top}{\cset{\kappa}}})}{\cempty}).\,\unboxc{z}{\cempty}{\ev{member}\,x\,(\ev{refl}\,\mathit{Bad})\,0} \\
\qquad \ann{\ev{elem}\,\cset{z}\,\cset{z}}\ \ann{\ev{elem}\,\cempty\,\cset{x}} .
\end{array}
\]
Its typing is decided (\lean{C3\_typed}), and so is $\ev{member}\,x\,(\ev{refl}\,\mathit{Bad})\,0 : \cset{\kappa} \le \cempty$ under both binders (\lean{C3\_badBounds}). The unboxing takes $\kappa$ off the boxed value and charges $\cempty$. Over a typed store no atom $a$ has a type $\capt{\mathit{Bad}}{D}$, since $\ev{member}\,a\,(\ev{refl}\,\mathit{Bad})\,0$ would be closed evidence that Corollary~\ref{cor:target-star} rules out. A call of the outer lambda from a typed state would hand its body such an atom, so the body never runs. The Lean states the corollary at such a binder as \lean{C3\_no\_store}, by hand, and the rest is report~1's argument for bad bounds on paper, with the roots of capture sets in place of head forms.
\end{example}

\subsection{The use-set lemma and the rest of the metatheory}\label{sec:target-meta}

The use-set half of preservation keeps the shape of Lemma~\ref{lem:mini-step} and gains a renaming $\pi$ of the de Bruijn binders of one store into the next, the identity with names.

\begin{lemma}[The use-set half of preservation, \leanin{step\_uses}{\treeK}]\label{lem:target-step}
Let $\statej{\mathit{st}}{T}$, let $\storej{\sigma}{\Gamma}$ for the store $\sigma$ of $\mathit{st}$, and let $\mathit{st} \step \mathit{st}'$, where $\sigma'$ is the store of $\mathit{st}'$. Then $\sextr{\pi}{\sigma}{\sigma'}$ for some renaming $\pi$, and $\caple{\Gamma'}{\uses{\mathit{st}'}}{\pi(\uses{\mathit{st}})}$ for every $\Gamma'$ with $\storej{\sigma'}{\Gamma'}$.
\end{lemma}

Table~\ref{tab:target-steps} gives the twelve cases. With the lemma, Theorems~\ref{thm:mini-pred} and~\ref{thm:mini-effect} hold for the full calculus with $\pi$ added, as \lean{capture\_prediction} and \lean{effect\_safety} in \treeK. Report~1's other theorems keep their form. A state is a cast redex when its next step only moves a cast frame, by \rn{E-CastPush}, \rn{E-CastVal} or \rn{E-CastAtom} (\lean{State.CastRedex}). Such a step erases to no runtime step.

\begin{table}[t]
\caption{The twelve steps of the \treeK machine and the case of Lemma~\ref{lem:target-step} that covers each. Items are those of Theorem~\ref{thm:target-canon}. A chain is the head form of the chain of casts of an atom. Four cases read one through \lean{Store.Typed.formsTyped}, which report~1's canonical forms supply.}
\label{tab:target-steps}
\footnotesize
\begin{tabular}{@{}p{0.25\linewidth}p{0.71\linewidth}@{}}
\toprule
Step & Case \\
\midrule
\mbox{\rn{E-Let}, \rn{E-CastPush},}\newline \mbox{\rn{E-CastVal}, \rn{E-CastAtom}} & the two use sets have the same atoms \\
\rn{E-Alloc} & the frame's avoidance evidence at the new transparent binder, then item 1 \\
\rn{E-Rename} & the frame's avoidance evidence at the atom, then item 1 \\
\rn{E-App} & the closing evidence of the stored closure at the argument, then item 1 \\
\rn{E-AppCastRefl} & as \rn{E-App}, with the chain $\form{id}$ or $\form{eqv}\,\varphi$, which makes the atom's function shape the closure's \\
\rn{E-AppCast} & as \rn{E-App} at $\ev{cast}\,b\,d$, with the chain $\form{pi}\,d\,c$, which types $d$ \\
\rn{E-Proj} & the closing evidence of the stored field at the object's variable, then item 1 \\
\rn{E-UnboxRefl} & item 3 at $b$, typed by the chain $\form{id}$ or $\form{eqv}\,\varphi$, then item 1 at the charge \\
\rn{E-UnboxCast} & item 3 at $\ev{cast}\,b\,d$, typed by the chain $\form{boxed}\,d$, then item 1 at the charge \\
\bottomrule
\end{tabular}
\end{table}

\begin{theorem}[Type safety and erasure, \lean{preservation'}, \lean{progress}, \lean{erase\_step} and \lean{erase\_reflect'} in \treeK]\label{thm:target-meta}
\leavevmode
\begin{enumerate}
\item If $\statej{\mathit{st}}{T}$ and $\mathit{st} \step \mathit{st}'$, then $\statej{\mathit{st}'}{\pi(T)}$ for some renaming $\pi$.
\item If $\statej{\mathit{st}}{T}$, then $\mathit{st}$ is final or $\mathit{st} \step \mathit{st}'$ for some $\mathit{st}'$.
\item If $\mathit{st} \step \mathit{st}'$, then $\mathit{st}$ is a cast redex and $\erase{\mathit{st}} = \erase{\mathit{st}'}$, or $\erase{\mathit{st}} \step \erase{\mathit{st}'}$ in the runtime.
\item If $\storej{\sigma}{\Gamma}$ for the store $\sigma$ of $\mathit{st}$, $\tmj{\Gamma}{t}{T}$ for its term $t$ and some $T$, and $\erase{\mathit{st}} \step r$ in the runtime, then $\mathit{st} \steps \mathit{st}'$ and $\erase{\mathit{st}'} = r$ for some $\mathit{st}'$.
\end{enumerate}
\end{theorem}

\subsection{Boxes}\label{sec:target-boxes}

Boxing hides a capture set and unboxing recovers it (Section~\ref{sec:cc-later}), and \treeK keeps both explicit. Boxing is a value, $\ev{box}\,a : \capt{(\boxt{T})}{\cempty}$ (\rn{V-Box}), so a container of capabilities can be pure, and a store location that holds a box has no roots. Unboxing is a term, $\unboxc{a}{U}{f}$, with the charge $f : C \le U$ of the boxed set $C$ against the declared use set $U$ (\rn{T-Unbox}). The capture set $D$ of the box itself plays no part. \rn{E-UnboxRefl} hands back the stored atom at the chain $\form{id}$ or $\form{eqv}\,\varphi$, also at a bare variable. \rn{E-UnboxCast} hands it back under $d$ at the chain $\form{boxed}\,d$, as \rn{E-AppCast} hands a closure its argument under the domain coercion. Progress and the backward erasure simulation use the next theorem, and the use-set lemma uses its typed chain.

\begin{theorem}[Closed box inversion, \leanin{closed\_box\_inversion}{\treeK}]\label{thm:target-box}
Let $\storej{\sigma}{\Gamma}$ and $\atomj{\Gamma}{a}{\capt{(\boxt{T})}{D}}$. Then $\sigma(\rt{a}) = \ev{box}\,b$ for some atom $b$, and $\chn{\sigma}{a}{n}{\Phi}$ for some $n$ and a head form $\Phi$ that is $\form{id}$, $\form{eqv}\,\varphi$ or $\form{boxed}\,d$.
\end{theorem}

C7 is a pure container of two boxed capabilities (\lean{C7\_client}, decided). Its object type is
\[
\mathit{Pair} = \obj{\tel_1 \mathbin{+\!\!+} \tel_2} \quad\text{with}\quad \tel_i = [\,\phas{\ell_i},\ \ple{\sel{\self}{\ell_i}}{\boxt{\CapT{\kappa_i}}},\ \plec{\cset{\sel{\self}{\ell_i}}}{\cempty}\,] .
\]
Both fields declare $\cempty$, so $o : \capt{\mathit{Pair}}{\cempty}$ is pure. The client binds $\ell_1$ of $o$ to $y : \capt{(\boxt{\CapT{\kappa_1}})}{\cempty}$ and opens it with $\unboxc{y}{\cset{\kappa_1}}{\ev{refl}\,\cset{\kappa_1}}$ in a let that declares $\cset{\kappa_1}$. With the declared sets $\cempty$ and the charge $\ev{elem}\,\cset{\kappa_1}\,\cempty$ the checker refuses the client (\lean{C7\_rejected}, decided).

Effect safety never flags an unboxing. The unboxing inspects the variable $\rt{a}$, which holds a box, so $\cset{\rt{a}}$ has no roots. The next theorem speaks of the atom it hands back.

\begin{theorem}[Effect safety at an unboxing, \leanin{effect\_safety\_unbox}{\treeK}]\label{thm:target-unbox}
Let $\statej{\mathit{st}}{T}$, let $\storej{\sigma}{\Gamma}$ for the store $\sigma$ of $\mathit{st}$, and let $\kappa \notin \roots{\Gamma}{\uses{\mathit{st}}}$. Let $\mathit{st} \steps \mst{\sigma'}{K'}{\unboxc{a}{U}{f}}$ with $\sigma'(\rt{a}) = \ev{box}\,b$, and let $\storej{\sigma'}{\Gamma'}$. Then $\sextr{\pi}{\sigma}{\sigma'}$ and $\pi(\kappa) \notin \roots{\Gamma'}{\cset{\rt{b}}}$ for some renaming $\pi$.
\end{theorem}

Here $\rt{b}$ is report~1's root of the atom $b$, a term variable, and the conclusion is about the roots of the capture set $\cset{\rt{b}}$. Its proof takes the unboxing step from progress and applies capture prediction to the longer run.

\subsection{What is not proved}\label{sec:target-open}

Three examples in \lean{Captures/Frontend/StepFC.lean} run an unboxing step on states that no Lean statement types (decided). No example runs one from a typed state, and no Lean statement instantiates Theorem~\ref{thm:target-unbox}. \lean{C7\_rejected} rules out one evidence term, and Corollary~\ref{cor:target-star} rules out all of them only over a typed store.
